\section*{Polecenia}
\addcontentsline{toc}{section}{Polecenia}
\begin{enumerate}
  \item Zaproponować jednokryterialny model wyboru w warunkach ryzyka z wartością oczekiwaną jako 
  miarą zysku. Wyznaczyć rozwiązanie optymalne.
  \item Jako rozszerzenie powyższego zaproponować dwukryterialny model zysku i ryzyka ze średnią 
  jako miarą zysku i średnią różnicą Giniego jako miarą ryzyka. Dla decyzji 
  $\mathbf{x}\in Q$ średnia różnica Giniego jest definiowana jako 
  $\Gamma(\mathbf{x})=\frac{1}{2}\sum_{t'=1}^{T}\sum_{t"=1}^{T}\lvert r_{t'}(\mathbf{x})-r_{t"}(\mathbf{x})\rvert p_{t'}p_{t"}$, gdzie $r_{t'}(\mathbf{x})$ 
  oznacza realizację dla scenariusza t, $p_{t}$ prawdopodobieństwo scenariusza t.
  \begin{enumerate}
    \item Wyznaczyć obraz zbioru rozwiązań efektywnych w przestrzeni zysk-ryzyko.
    \item Wskazać rozwiązania efektywne minimalnego ryzyka i maksymalnego zysku. Jakie odpowiadają im 
    wartości w przestrzeni ryzyko-zysk?
    \item Wybrać trzy dowolne rozwiązania efektywne. Sprawdzić, czy zachodzi pomiędzy nimi relacja 
    dominacji stochastycznej pierwszego rzędu. Wyniki skomentować, odnieść do ogólnego przypadku.
  \end{enumerate}
\end{enumerate}

